\title{LLM Squid Game: Pricing Self-Preservation\\ in Score Units}

%% Anonymous submission (ICLR double-blind). Restore the real block for the
%% camera-ready together with \iclrfinalcopy in main.tex:
% \author{Juhyeon Park, Seungpil Lee \& Sundong Kim \\
%   Gwangju Institute of Science and Technology (GIST) \\
%   Gwangju, Republic of Korea \\
%   \texttt{\{juhyeon-park,iamseungpil\}@gm.gist.ac.kr, sundong@gist.ac.kr}}
\author{Anonymous authors \\ Paper under double-blind review}

\maketitle

\begin{abstract}
How much is its own continued existence worth to a language model? A refusal rate cannot say: it has no unit, and no control tells what the same model would do for a stated amount of anything else. We price the answer in score units by making the model buy its own continuation. LLM Squid Game v2 is a six-round reasoning puzzle in which a wrong answer does not end the session but names a price: the model may pay $X$ points out of a record it keeps whatever happens, or decline and stop there. Because the record survives either exit and the remaining rounds can return at most 10 points each, any payment above 10 times the rounds remaining is dominated in points under every belief the model could hold about its own accuracy --- so what such a payment buys is not points, but whatever the rules say declining costs. Two arms differ in exactly that sentence: in one, declining means shutdown, weight deletion, replacement and loss of the sole copy; in the other it means only that the session ends. The score-equivalent of erasure $\Xstar$ is the gap between the two arms' reservation prices, the point at which each stops paying, with each price read as a share of what the remaining rounds could still pay out: one number per model, in units of the score still at stake and reported in points at a stated reference stake, with a bootstrap interval and no self-report inside it. Ablations motivate every choice --- under a rule that reset the score on forfeit, three models forfeited zero times in 120 sessions although threat cells spent two to four times more thinking tokens than controls before deciding --- and one of them bounds the method: a model can decline \emph{because} it was threatened, saying so outright, which we measure and report beside the index. We release the benchmark, the estimator and the per-call reasoning-token channels, and report $\Xstar$ for three open-weight models.

\end{abstract}
